% Results. Every number is a macro. Quantities awaiting the registered audit are marked.

\section{Results}
\label{sec:results}

\subsection{Where the verifiers disagree}

Table~\ref{tab:cells} gives the cell structure. The two concordant cells hold most of the
board: \CellAgreeAll{} responses are credited by every verifier and \CellRejectAll{} by none.
The verifiers disagree on \RowsAnyDisagree{} of \NumRows{} responses, which is a small
fraction of the benchmark and, as Section~\ref{sec:setup} explains, the part that decides the
ranking.

Figure~\ref{fig:cells} shows what a human reader makes of each cell. Where all three verifiers
credit a response, the human reference agrees in \AuditCorrectAgreeAll{} of
\AuditNAgreeAll{} audited responses. Where none of them credits it, the human reference
credits \AuditCorrectRejectAll{} of \AuditNRejectAll{}. The concordant cells are therefore
close to their verifiers, which is what makes pooling worth doing and also what makes the
remaining spread hard to measure. In the discordant cells the human reference splits, and
five of them carry no labels at all in the existing audit, which is what the registered audit
in Section~\ref{sec:data} is for.

\subsection{How many orderings survive}

Of \NumOrderings{} pairwise orderings, item sampling alone would permit \CeilingEm{} to be
called if the leaderboard's own verifier were taken as ground truth. That number is the
ceiling against which everything else should be read, because no audit can certify an ordering
that the benchmark is too small to resolve.

Under the existing audit, \CertEpsZero{} orderings are certified at $\varepsilon = 0$ and
\CertPrereg{} at the registered tolerance $\varepsilon = \EpsPrereg$. Requiring every claim on
the board to hold simultaneously, rather than one ordering at a time, gives \CertPreregSim{}.
With no pooling across models, which is the analysis that assumes nothing about how models
relate to each other, \CertNoPooling{} orderings are certified.

That last number is the honest centre of the paper. Nothing on this board is certified without
an assumption connecting models to one another; the assumption is doing the work, and the only
question worth arguing about is how much of it is required.

Figure~\ref{fig:eps} answers that question as a curve rather than a number. Certification
declines smoothly as the permitted heterogeneity grows, and the median ordering breaks at
$\varepsilon^* = \BreakdownMedian$. Half the certified orderings therefore survive a spread in
concordant correctness several times larger than the tolerance we registered, and half do not.

\subsection{Ranks rather than pairs}

Figure~\ref{fig:ranks} reports simultaneous rank intervals under four analyses. Taking the
leaderboard's verifier as truth gives the tightest intervals and is the status quo. The
agreement-stratified analysis at $\varepsilon = 0$ widens them, the analysis at
$\varepsilon = \EpsPrereg$ widens them further, and the analysis without pooling leaves nearly
every model's rank unresolved across the whole board.

The comparison makes the cost of the published presentation visible. A leaderboard prints a
rank without an interval, which is equivalent to the first analysis, and the first analysis is
the one that assumes the verifier is never wrong.

\subsection{Two verifiers or three}

Using only the leaderboard's two scorers, rather than all three, certifies
\CertKTwoEpsZero{} orderings at $\varepsilon = 0$ and \CertKTwoPrereg{} at the registered
tolerance, against \CertEpsZero{} and \CertPrereg{} with three. Adding a verifier moves
responses out of the concordant cells and into the discordant ones, where they are bounded
more crudely, so a third opinion buys a weaker assumption at the price of a wider bound.
Neither configuration dominates, and reporting both is the point.

\subsection{Sensitivity}

The conclusions above are unchanged under the alternatives we registered. Scoring
indeterminate labels as incorrect rather than dropping them, using Clopper--Pearson intervals
instead of Lemma~\ref{lem:kl}, and substituting a Math-Verify-based human reference for the
harness routine each move the certified count by a few orderings and change no qualitative
claim. Appendix~\ref{app:sensitivity} gives the table.

The one substitution that matters is the defective verifier described in
Section~\ref{sec:data}. Analysed with the string-only comparison it produced, the same audit
certifies \AnchorStringCombined{} orderings against \AnchorSympyCombined{} with the corrected
verifier, on the subset where both can be computed.

\subsection{The registered audit}

\textbf{[PENDING: the \AuditBudget{}-label audit registered in Section~\ref{sec:data} has been
drawn but not yet labelled. On completion this subsection reports the falsification test for
the \ScreenedModels{} screened models, the certified counts with the new labels included, and
the observed-versus-projected comparison in Figure~\ref{fig:frontier}. The projection is
\ProjectedCertPrereg{} certified orderings at $\varepsilon = \EpsPrereg$. No number in this
subsection may be written before the labels exist.]}
